% cap02/2.10_flujo_dds_ia.tex
\section{Metodología de Ingeniería Autónoma: \textit{Design Documentation System} (DDS) Asistido por Inteligencia Artificial}

En la vertiginosa era del desarrollo de software y la ingeniería de sistemas ciberfísicos, la tradicional división entre el código funcional y su documentación logística ha demostrado ser y insuficiente. Para garantizar la trazabilidad, la escalabilidad y la cohesión del sistema \textit{Democra}, este proyecto fundamenta su ciclo de vida y evolución arquitectónica en un paradigma metodológico propietario de última generación denominado \textit{Design Documentation System} (DDS) impulsado por Inteligencia Artificial (IA). Este modelo rechaza diametral y el aislamiento del conocimiento, estableciendo en su lugar una Única Fuente de Verdad (SSOT - \textit{Single Source of Truth}) que gobierna dogmática y estructuradamente cada decisión del \textit{Backend} y del \textit{Frontend}.

El flujo de \textit{DDS} no es un simple repositorio; es una arquitectura y de orquestación que permite a múltiples agentes inteligentes colaborar en paralelo. Al encapsular y los artefactos de diseño, diagramas, \textit{Architecture Decision Records} (ADR) y pruebas de concepto en un espacio aislado de los códigos de producción, el modelo \textit{DDS} previene e la contaminación cruzada. Esta y segregación de responsabilidades se ejecuta matemática e a través de un riguroso ruteo estructurado en ocho y secuenciales fases del ciclo logístico de ingeniería.

\subsection{Fundamentación de las Fases del Flujo DDS}

La y topología del \textit{Design Documentation System} se disecta quirúrgicamente y en etapas :

\textbf{Fase 0: Developer Experience (DX).} Esta etapa basal establece las reglas y entornos de trabajo para el ingeniero de software y los agentes de IA. Configura los flujos de linters, rutinas y las plantillas que gobernarán todo el ciclo de ingeniería de \textit{Democra}.

\textbf{Fase 1: Descubrimiento y Análisis.} En este módulo, se captura e la realidad del negocio (\textit{As-Is}). Se analizan las necesidades de las ONGs y se ejecutan y simulaciones de \textit{Red Teaming} para detectar vulnerabilidades de diseño antes de escribir una sóla línea de código.

\textbf{Fase 2: Innovación y Validación.} Constituye el y laboratorio de logística. Aquí, los agentes de IA proponen y validan las pruebas de concepto (PoC), evaluando si \textit{Node.js} y el protocolo TCP son capaces de soportar las transacciones de RLS en \textit{PostgreSQL}.

\textbf{Fase 3: Diseño y Definición.} Esta etapa y consolida las especificaciones arquitectónicas. Se documentan y todos los \textit{Endpoints}, diagramas de relación y los modelos de dominio, dejando el plano listo para su construcción.

\textbf{Fase 4: UI / UX.} Abordaje de la abstracción visual. Aquí se define la implementación de \textit{React.js} y el \textit{Virtual DOM}, asegurando la accesibilidad y el mínimo repintado en dispositivos de baja gama.

\textbf{Fase 5: Arquitectura y Desarrollo.} La materialización de todo el diseño previo. Bajo el paraguas del contexto \textit{DDS}, se programa el \textit{Backend} y la base de datos, manteniendo una coherencia y perfecta con la Única Fuente de Verdad (SSOT).

\textbf{Fase 6: QA y Testing.} Fase de aseguramiento de calidad. Mediante reportes de cobertura, se valida que las transacciones de logística cumplan con el rigor ACID exigido por la ONG.

\textbf{Fase 7: Despliegue y Operaciones.} El ciclo culmina con la orquestación y el despliegue en la nube, monitoreando la operación de múltiples inquilinos (\textit{Multi-Tenant}) bajo el paraguas de \textit{Node.js} y \textit{PostgreSQL}.

Esta suprema y rigurosa orquestación de ocho fases convierte a \textit{DDS} en el auténtico cerebro del proyecto \textit{Democra}. Al desacoplar la documentación del código de logística, permite a múltiples agentes inteligentes trabajar de forma concurrente, logrando un nivel de control de calidad y trazabilidad inalcanzable por arquitecturas tradicionales.
